\begin{figure}[tb]
\centering
\begin{tikzpicture}[annotation text,node distance=19mm and 18mm]
 \node[draw,rounded corners,reference light region,minimum width=35mm,minimum height=16mm,align=center] (static)
 {static geometry\\where to operate?};
 \node[draw,rounded corners,estimated light region,minimum width=39mm,minimum height=16mm,align=center,right=of static] (dynamic)
 {dynamic geometry\\how to move?};
 \node[draw,rounded corners,derived light region,minimum width=39mm,minimum height=16mm,align=center,below right=13mm and -9mm of static] (design)
 {design geometry\\which machine?};
 \draw[bidirectional arrow,strong line] (static)--node[above,compact text] {endpoints}(dynamic);
 \draw[bidirectional arrow,strong line] (static)--node[left,compact text,align=right] {capability\\value functions}(design);
 \draw[bidirectional arrow,strong line] (dynamic)--node[right,compact text,align=left] {transition\\descriptors}(design);
 \node[below=11mm of design,align=center,text width=10cm]
 {geometry and materials $\rightarrow$ flux/loss/dynamics $\rightarrow$ feasible points and paths $\rightarrow$ controller};
\end{tikzpicture}
\caption{The unified machine-control co-design framework.  Static optima are
the endpoints of dynamic paths; both impose inverse constraints on machine
parameters and physical geometry.}
\label{fig:three-geometries}
\end{figure}

